\begin{example}[A column indexes local vocabulary banks, not a common token]
\label{ex:algorithm-banks-20260918}
Take $N=15$, $G=\emptyset$, and a vocabulary whose first three tokens in
the public order are $a\prec b\prec c$ (further tokens may follow).
Suppose the local vocabulary banks $\TokenBank{i}$ from
\AlgoBlockRef{preprocess} (\cref{app:def:preprocessing}) and the fixed draft are
\[
 \TokenBank{1}=\{a,c\},\quad \DraftToken{1}=a,
 \qquad
 \TokenBank{i}=\{b\},\quad \DraftToken{i}=b
 \quad (i\in[15]\setminus\{1\}).
\]
This specifies the screen inputs, not a target distribution or an oracle.
In particular, $\BankSize{1}=2$, $\BankSize{i}=1$ for $i\ne1$, and
$\MaxBankSize=2$.  The first token outside each bank is
$\TailRepresentative{1}=b$ and $\TailRepresentative{i}=a$ for $i\ne1$.
Substituting into \eqref{app:eq:packed-column} gives
\[
\begin{array}{c|ccc}
 \text{position}&\kappa=1&\kappa=2&\kappa=\TailColumn\\ \hline
 1& a&c&b\\
 i\ne1&b&\DraftToken{i}=b&a
\end{array}
 \qquad
 \ScreenColumnSet=\{1,2,\TailColumn\},\quad \ScreenColumnCount=3.
\]
Thus column $1$ assigns $a$ at position $1$ but $b$ at position $4$.
Column $2$ uses the second bank token at position $1$ and draft padding
at position $4$, whose bank has no second token.  The tail column uses
the representative outside each bank.  The same three columns are
reused for every source color and readout chunk.
\end{example}
